% Problem record op_428e6c37ba03149a. This TeX file is derived from
% database/problems_json/op_428e6c37ba03149a.json by scripts/sync-tex.mjs; edit the JSON record
% and rerun the script rather than editing this file.
\section{Optimal dimension-free copy complexity of shadow tomography}
\label{sec:op_428e6c37ba03149a}

\paragraph{Problem.}
Can list-dependent shadow tomography always achieve dimension-free copy complexity $O(\log M/\varepsilon^2)$?

Let $f_1,\ldots,f_M:\mathcal X\to[0,1]$ be known functions and let $X_1,\ldots,X_N$ be independent samples from an unknown distribution. Estimate every expectation with its empirical mean. Hoeffding's inequality and a union bound give

\begin{equation}
\Pr\!\left[\max_j|\widehat\mu_j-\mathbb E f_j(X)|>\varepsilon\right]
\le 2M e^{-2N\varepsilon^2}.
\label{eq:428e-1}
\end{equation}

Thus $N\ge\ln(6M)/(2\varepsilon^2)$ suffices for simultaneous success probability at least $2/3$, independently of the size of $\mathcal X$. This is the classical benchmark; one observed sample can be evaluated against every known function.

The quantum task fixes $M\ge2$, an arbitrary finite dimension $d$, and a known list of effects

\begin{equation}
0\le E_1,\ldots,E_M\le I_d.
\label{eq:428e-2}
\end{equation}

An unknown state $\rho$ is supplied as independent copies. Question: Does a universal constant $C$ exist such that, for every list and every $0<\varepsilon<1/4$, a measurement and estimator using

\begin{equation}
N\le\left\lceil C\frac{\log M}{\varepsilon^2}\right\rceil
\label{eq:428e-3}
\end{equation}

can produce estimates satisfying

\begin{equation}
\inf_{\rho}\Pr_\rho\!\left[
\max_{1\le j\le M}|\widehat\mu_j-\operatorname{Tr}(\rho E_j)|\le\varepsilon
\right]\ge\frac23?
\label{eq:428e-4}
\end{equation}

The list is known before measurement. Arbitrary collective measurements on all copies are allowed, and no computational-efficiency requirement is imposed. This is the list-dependent shadow-tomography task, not measurement-independent classical shadows or adaptively chosen online queries. \sourcecite{ref:428e-aaronson18}{Aaronson18}, \sourcecite{ref:428e-jeronimo26}{Jeronimo26}

If the effects commute, measuring their common eigenbasis reduces the task to the classical empirical-estimation argument above. Therefore the noncommuting case is the only obstacle in the displayed comparison.

The displayed definitions, constraints, and target bounds are recorded in Eqs.~\eqref{eq:428e-1}, \eqref{eq:428e-2}, \eqref{eq:428e-3}, \eqref{eq:428e-4}.

\subsection*{Status}

Unsolved

\subsection*{Source}

The question is explicitly posed or retained as open in the cited primary literature \sourcecite{ref:428e-aaronson18}{Aaronson18}\sourcecite{ref:428e-jeronimo26}{Jeronimo26}.  The statement is rewritten here to make its hypotheses and success criterion self-contained.

\subsection*{Progress}

\begin{itemize}
  \item Aaronson's shadow-tomography work established a worst-case lower bound of order

\begin{equation}
\Omega\!\left(\frac{\min\{d^2,\log M\}}{\varepsilon^2}\right).
\label{eq:428e-5}
\end{equation}

The displayed definitions, constraints, and target bounds are recorded in Eqs.~\eqref{eq:428e-5}.

  \item For sufficiently large dimension, this matches the proposed target. It is not a lower bound of $\log M/\varepsilon^2$ for every fixed small $d$. \sourcecite{ref:428e-aaronson18}{Aaronson18}

  \item Chen, Li, and Liu proved the optimal logarithmic rate in a high-precision regime. Subsequent work by Pelecanos, Spilecki, and Wright expands the regime where that rate is achievable. These dimension-dependent precision conditions do not settle all $d$ and $\varepsilon$ simultaneously. \sourcecite{ref:428e-chen24}{Chen24}, \sourcecite{ref:428e-pelecanos25}{Pelecanos25}

  \item The critical update is Jeronimo, Huang, and Liu, arXiv:2608.06345v2, revised September 14, 2026. The revision reports

\begin{equation}
N=O\!\left(\frac{\log M\,\log(M/\delta)}{\varepsilon^2}\right).
\label{eq:428e-6}
\end{equation}

The displayed definitions, constraints, and target bounds are recorded in Eqs.~\eqref{eq:428e-6}.

  \item At failure probability $\delta=1/3$, this becomes $O(\varepsilon^{-2}\log^2 M)$. The earlier version's fourth-power logarithm is no longer the latest reported bound. The theorem also answers the older question of whether dimension-free polylogarithmic shadow tomography exists. \sourcecite{ref:428e-jeronimo26}{Jeronimo26}
\end{itemize}

\subsection*{References}

\begin{enumerate}[label={},leftmargin=4.8em,labelsep=0.5em]
  \footnotesize
  \item[\textup{[Aaronson18]}]\label{ref:428e-aaronson18}
  S. Aaronson, \emph{Shadow Tomography of Quantum States}, arXiv:1711.01053; STOC (2018). \href{https://arxiv.org/abs/1711.01053}{Paper}.

  \item[\textup{[Jeronimo26]}]\label{ref:428e-jeronimo26}
  F. Granha Jeronimo, Q. Huang, and L. Liu, \emph{Dimension-Free Polylogarithmic Quantum Shadow Tomography from Sequential Pretty-Good Measurements}, arXiv:2608.06345v2, revised September 14, 2026. The HTML manuscript uses the shorter title \emph{Dimension-Free Polylogarithmic Quantum Shadow Tomography}. See Definition 1.1, Theorem 1.2, and Section 1.2. \href{https://arxiv.org/abs/2608.06345v2}{Versioned record}; \href{https://arxiv.org/html/2608.06345v2}{versioned full text}.

  \item[\textup{[Chen24]}]\label{ref:428e-chen24}
  S. Chen, J. Li, and A. Liu, \emph{Optimal high-precision shadow estimation}, arXiv:2407.13874 (2024). \href{https://arxiv.org/abs/2407.13874}{Paper}.

  \item[\textup{[Pelecanos25]}]\label{ref:428e-pelecanos25}
  A. Pelecanos, J. Spilecki, and J. Wright, \emph{The debiased Keyl's algorithm: a new unbiased estimator for full state tomography}, arXiv:2510.07788 (2025). The improvement to the high-precision regime is also discussed in Section 1.2 of \sourcecite{ref:428e-jeronimo26}{Jeronimo26}. \href{https://arxiv.org/abs/2510.07788}{Paper}.
\end{enumerate}

\subsection*{Comment}

At constant failure probability, the latest reported bounds leave the worst-case gap

\begin{equation}
\Omega(\varepsilon^{-2}\log M)
\quad\text{versus}\quad
O(\varepsilon^{-2}\log^2 M).
\label{eq:428e-7}
\end{equation}

The result is currently available as a recent preprint. The open question is stated as an explicit optimization target, not attributed as a newly named conjecture to those authors.

The displayed definitions, constraints, and target bounds are recorded in Eqs.~\eqref{eq:428e-7}.

\subsection*{Field}

Quantum metrology

\subsection*{Topic}

Quantum estimation; One-shot and finite-blocklength bounds

\subsection*{ID}

\texttt{op\_428e6c37ba03149a}
